% scrartcl = moderne, deutsche Artikel-Klasse (Überschriften standardmäßig linksbündig)
\documentclass[a4paper, 12pt, titlepage, parskip=half, twoside, headinclude, footnotes=bftn]{scrartcl}

\usepackage[utf8]{inputenc}
\usepackage[ngerman]{babel} % Für deutsche Rechtschreibung und korrekte Silbentrennung 
\usepackage{amsmath}
\usepackage{setspace}
\setstretch{1.14} % Skaliert den Zeilenabstand exakt um den Faktor 1,08
\usepackage{booktabs}  % Für professionelle Tabellenlinien (\toprule, \midrule, \bottomrule)
\usepackage{xcolor}
\definecolor{darkblue}{rgb}{0.0, 0.0, 0.55}

% ==========================================
% GLOBALE METADATEN (Hier zentral anpassen!)
% ==========================================
\newcommand{\mytitle}{Berechnung des Trocknungsfortschritts aus den Messwerten relative Feuchte und Temperatur}
\newcommand{\myshorttitle}{Berechnung des Trocknungsfortschritts}
\newcommand{\myauthor}{Burkhard Pick}      % Autor(en) des Dokuments
\newcommand{\myversion}{Version 0.1}              % Aktuelle Version
\newcommand{\mystatus}{Entwurf}            % Status: z.B. Entwurf, Review, Final
\newcommand{\mydate}{\today}               %

% Filtert alle reinen Schrift-Informationsmeldungen aus dem Terminal
\usepackage{silence}
\WarningFilter{latex}{Font shape} 
\WarningFilter{latex}{Font Info}

% ==========================================
% MODERNE SCHRIFT FÜR PDFLATEX (Fira Sans)
% ==========================================
\usepackage[sfdefault]{FiraSans} %% sfdefault setzt sie als Hauptschrift
\usepackage[T1]{fontenc} % Wichtig für korrekte Silbentrennung bei PostScript-Fonts
\usepackage{FiraMono}   %% Passende Festbreitenschrift für Code

\usepackage[
    top=2.5cm,      % Rand oben
    bottom=3.0cm,   % Rand unten
    left=2.0cm,     % Rand links
    right=2.0cm,    % Rand rechts
    bindingoffset=0mm, % ZUSÄTZLICHER Platz für die Bindung (wird innen hinzugefügt)
    footskip=1.5cm%,  % Abstand vom Text zur Fußzeile
    %headinclude=true % Kopfzeile als Teil des Textbereichs (des sogenannten Satzspiegels) gewertet 
]{geometry} % Standardmäßig sind LaTeX-Ränder sehr breit. Hier sind sie kompakt eingestellt

% KOPF- UND FUSZEILEN MIT LINIEN (scrlayer-scrpage)
\usepackage[
    headsepline=0.4pt, % Linie unter der Kopfzeile (Dicke: 0.4 Punkte)
    footsepline=0.4pt, % Linie über der Fußzeile (Dicke: 0.4 Punkte)
    automark           % Aktiviert das automatische Aktualisieren der Überschriften
]{scrlayer-scrpage}
\clearpairofpagestyles % Löscht alle Standard-Kopf- und Fußzeilen

\usepackage{lastpage}
% --- KOPFZEILE (Mit Section-Überschrift) ---
% Da du twoside nutzt, wird die Überschrift außen platziert:
\lehead{Berechnung des Trocknungsfortschritts} % Überschrift links auf geraden Seiten
\rohead{Berechnung des Trocknungsfortschritts} % Überschrift rechts auf ungeraden Seiten

% Seitennummer in der Fußzeile außen platzieren (zweiseitiges Layout):
\lefoot{Seite~\pagemark{} von~\pageref{LastPage}} % Links auf geraden Seiten
\rofoot{Seite~\pagemark{} von~\pageref{LastPage}} % Rechts auf ungeraden Seiten

% Text auf der Innenseite (bei zweiseitigem Druck):
\refoot{\today} % rechts auf geraden Seiten (innen)
\lofoot{\today} % links auf ungeraden Seiten (innen)
\cfoot{\myversion}

\usepackage{minted} % Das mächtige Paket für echtes Syntax-Highlighting

% ==========================================
% LITERATURVERZEICHNIS MIT BIBLATEX
% ==========================================
\usepackage[
    backend=biber,       % Nutzt Biber als modernen Sortierer
    style=numeric-comp,  % Zitierstil: Nummern z.B., [2] (Kompakt bei mehreren)
    sorting=nty,         % Sortierung: Name, Titel, Jahr
    autolang=other       % Richtige Trennung für fremdsprachige Quellen
]{biblatex}
\addbibresource{literatur.bib} % Lädt deine erstellte .bib-Datei

\usemintedstyle{xcode} % Markdown-Theme z.B. github, monokai, vs, xcode

\title{\mytitle}
\author{\myauthor}
\date{\mydate}

\usepackage[
    colorlinks=true, % Aktiviert farbigen Text statt Rahmen
    linkcolor=black, % Interne Links (z.B. Inhaltsverzeichnis) bleiben schwarz
    citecolor=black, % Literaturverweise bleiben schwarz
    urlcolor=blue,    % Web-Links (URLs) werden blau eingefärbt
    psdextra
]{hyperref}

\begin{document}
\raggedbottom % Schaltet den Blocksatz für die vertikale Seite ab.

\begin{titlepage}
\maketitle

% Enter abstract:
\end{titlepage}

\tableofcontents % Automatisches Inhaltsverzeichnis

\section{Berechnung der absoluten Feuchte in \texorpdfstring{$\mathrm{g} / \mathrm{kg}$}{g / kg}} 

%\section{Berechnung der absoluten Feuchte in \textorpdfstring{$\text{g} / \text{kg}$}{g / kg}} 


\subsection{Messwerte in der Dämmschichtrocknung}
Berechnung der absoluten Feuchte aus gemessenen Werten:

\begin{itemize}
\item \textbf{relative Feuchtigkeit:} rF in Prozent im Wertebereich $[0 .. 100]$
\item \textbf{Temperatur:} T in °C: 
\item calcValue: absolute Feuchte aF in $\frac{\mathrm{g}}{\mathrm{kg}}$
\end{itemize}

In der POLYGON App und in der Apama-Regel in der CoT wird die Berechnug der absoluten Feuchte in $\frac{\mathrm{g}}{\mathrm{kg}}$ wie folgt durchgeführt:

\begin{minted}[
    bgcolor=gray!10,   % Grauer Hintergrundkasten
    fontsize=\small,   % Angenehme Schriftgröße
    breaklines         % Automatischer Zeilenumbruch
]{java}  
    const calcValue = ((((rF / 100) * (6.112 * Math.pow(2.7182881, (17.62 * T) / (243.12 + T)) * 100)) / 461.5 / (T + 273)) * 1000) / 1.2
\end{minted}

\subsection{Herleitung und Analyse der Formel}
Implementierte Formel:
\begin{equation}
    \text{aF}  
    = \frac{\frac{\frac{\text{rF}}{100} \cdot 6{,}112 \cdot 2{,}7182881^{\frac{17{,}62 \cdot \vartheta}{243{,}12 + \vartheta}} \cdot 100}
                  {461{,}5 \cdot (\vartheta + 273)} \cdot 1000}
            {1{,}2} 
\end{equation}

Schrittweise Vereinfachung:
\begin{equation}
    \text{aF}  
    = \frac{\text{rF} \cdot 6{,}112 \cdot e^{\frac{17{,}62 \cdot \vartheta}{243{,}12 + \vartheta}} }
            {461{,}5 \cdot (\vartheta + 273)}
        \cdot \frac{1000}{1,2}
\end{equation}

Abbilden auf die  \textbf{Tetens-Gleichung} bzw. \textbf{Magnus Gleichung} und anpassen auf die Einheit $\frac{\text{g}}{\text{kg}}$:
\begin{equation}
    \text{aF}  
    = \frac{\text{rF} \cdot 611{,}2\, \frac{\text{J}}{\text{m}^3} \cdot e^{\dfrac{17{,}62 \cdot \vartheta}{\vartheta + 243{,}12}}}
    {461{,}52 \frac{\text{J}}{\text{kg} \cdot \text{K}} \cdot (\vartheta + 273{,}15)\, \text{K}}
    \cdot \frac{10}{1{,}204\, \frac{\text{kg}}{\text{m}^3}}
\end{equation}
mit 
\begin{itemize}
    \item $\vartheta$ : Termodynamische Temperatur in \textdegree C mit $T = (273{,}15 +  \vartheta)\, \text{K}$
    \item Faktor 10: 
    \begin{itemize}
        \item Umrechnung Prozent im Wertebereich $[0 \ldots 100]$ in den Wertebereich $[0 \ldots 1]$
        \item Umrechung der Masse von kg in g 
    \end{itemize}
    \item Die Dichte trockener Luft bei bei Standardbedingungen (20 °C auf Meereshöhe: $101{,}3, \text{Pa}$) beträgt $1,204 \frac{\text{kg}}{\text{m}^3}$. 
\end{itemize}

Die Formel berechnet die absolute Feuchte nicht wie üblich in Gramm pro Kubikmeter ($\mathrm{g/m}^{\mathrm{3}}$), sondern im Mischungsverhältnis Gramm pro Kilogramm trockene Luft ($\mathrm{g/kg}$). Um von der Masse pro Volumen ($\mathrm{m}^{\mathrm{3}}$) auf die Masse pro Luftmasse ($\mathrm{kg}$) zu kommen, muss durch die Dichte der Luft geteilt werden.

Alle Einheiten können gekürzt werden. Es ergibt sich ein Verhältnis der Masse von Wasserdampf zur Masse von trockener Luft. Damit kann die Einheit  $\frac{\text{g}}{\text{kg}}$ gewählt werden.

Damit ist diese Formel gültig für einen Druck von $101{,}3\, \text{Pa}$. Der Druck in der Prozessluft der Dämmschichttrockung wird davon abweichen. Die Größe der Abweichung und deren Einfluss sind zu betrachten.


\section{Exkurs: Tetens-Gleichung bzw. Magnus Gleichung}

Der Partialdruck eines Gases $e$ innerhalb eines Gasgemisches kann durch Umstellung der  \href{https://de.wikipedia.org/wiki/Thermische_Zustandsgleichung_idealer_Gase}{allgemeinen Gasgleichung} 
\begin{equation}
    p \cdot V = n \cdot R_{m} \cdot T  \rightarrow   p = \frac{n}{V} \cdot R_m \cdot T
\end{equation}
mit folgender Formel näherungsweise berechnet werden:
\begin{equation}
    e = \rho \cdot R_m \cdot T
\end{equation}
mit
\begin{itemize}
    \item  $p$ : Druck $\text{Pa} = \frac{\text{N}}{\text{m}^2} = \frac{J}{\text{m} \cdot \text{m}^2} = \frac{J}{\text{m}^3}$ 
    \item  $V$: Volumen in $\text{m}^3$ 
    \item  $n$: Stoffmenge in Mol 
    \item  $R_{m}$: Gaskonstande 8,31446261815324 $\frac{\mathrm{J}}{\mathrm{mol} \ \mathrm{K}}$ 
    \item  $T$: Termodynamische Temperatur in Kelvin 
    \item $\vartheta$: Termodynamische Temperatur in °C mit $T = 273{,}15 +  \vartheta$
    \item $\rho$ : Dichte in $\frac{\text{kg}}{\text{m}^3}$
\end{itemize} 
Für die Dampdichte $\rho_v$ folgt:
\begin{equation}
    \rho_v = \frac{p}{R_v \cdot T}
\end{equation}
Mit der Gaskonstante für Wasserdampf: $R = 461,52\,\frac{\text{J}}{\text{kg}\,\text{K}}$ und Sättigungsdampfdruck $\rho(T)$ folgt für den Druck  $p$
\begin{equation}
    p = \varphi \cdot \rho(T)
\end{equation}

Damit kann die \textbf{absolute Feuchte} der Luft in $\text{g} / \text{m}^3$ mit der folgenden Formel berechnet werden:
\begin{equation}
    \rho_v(\varphi, T) = \frac{\varphi \cdot p_s(T)}{R_v \cdot T_K}
\end{equation}
mit 
\begin{itemize}
    \item $\varphi$ : relative Luftfeuchtigkeit im Wertebereich $[0 \dots  1]$
    \item $p_s(T)$ :  Sättigungsdampfdruck in Pa (*saturation*)
    \item $R_v = 461,52\, \frac{\text{J}}{\text{kg} \cdot \text{K}}$: Gaskonstante für Wasserdampf (*vapor*)
    \item $T_K = T_{{}^\circ\text{C}} + 273,15\, \text{K}$ : Temperatur in Kelvin
\end{itemize}

Der Sättigungsdampfdruck $p_s(T)$ kann mit der \href{https://www.chemie.de/lexikon/S%C3%A4ttigungsdampfdruck.html} {\textbf{Tetens-Gleichung} bzw. \textbf{Magnus Gleichung}} approximiert werden:
\begin{equation}
    p_s(\vartheta) = 611{,}2\, \text{Pa} \cdot e^{\dfrac{17{,}62 \cdot \vartheta}{\vartheta + 243{,}12^\circ\text{C}}}
\end{equation}
wobei $\vartheta$ die Temperatur in °C angegeben ist.
Es gilt: 
\begin{itemize}
    \item $-30\,{}^\circ\mathrm{C} \leq 70\,{}^\circ\mathrm{C}$
    \item $p_s(\vartheta = 0\,{}^\circ\mathrm{C}) = 6{,}11213\;\mathrm{hPa} = 611{,}213\;\mathrm{Pa}$
\end{itemize}

Die Konstanten variieren in der Literatur, z.\ B.:
\begin{equation}
    p_s(\vartheta) 
    = 
    611{,}153\, \text{Pa} \cdot e^{\frac{22{,}4433 \cdot \vartheta}{\vartheta + 272{,}186^\circ\text{C}}}
\end{equation}


\section{Absolute Feuchte mit Berücksichtigung des Drucks}
KI generierter Code für die absolute Feuchte unter Berücksichtigung des Drucks:
\begin{minted}[
    bgcolor=gray!10,   % Grauer Hintergrundkasten
    fontsize=\small,   % Angenehme Schriftgröße
    breaklines         % Automatischer Zeilenumbruch
]{python}
def absolute_feuchte(temp_c, rf_prozent, druck_pa):
    e_sat = 611.2 * round(2.718281828459045 ** ((17.62 * temp_c) / (temp_c + 243.12)), 10)
    # Using numpy exp or math exp for precision
    import math
    e_sat = 611.2 * math.exp((17.62 * temp_c) / (temp_c + 243.12))
    e = (rf_prozent / 100.0) * e_sat
    af = 622.0 * e / (druck_pa - e)
    return af

af_normal = absolute_feuchte(20, 85, 101325)
af_vakuum = absolute_feuchte(20, 85, 91325)
print(f"Normal: {af_normal:.2f} g/kg")
print(f"Vakuum: {af_vakuum:.2f} g/kg")
\end{minted}
Daraus ergibt sich:
\begin{equation}
    \text{aF}(\vartheta, rF, p) = 622 \cdot \frac{\frac{rF}{100} \cdot 611{,}2 \cdot e^{\frac{17{,}62 \cdot \vartheta}{\vartheta + 243{,}12}}}{p - \left(\frac{rF}{100} \cdot 611{,}2 \cdot e^{\frac{17{,}62 \cdot \vartheta}{\vartheta + 243{,}12}}\right)} \quad \left[\frac{\text{g}}{\text{kg}}\right]
\end{equation}
Mit  
\begin{align}
e_{\text{sat}}(\vartheta) &= 611{,}2 \cdot e^{\frac{17{,}62 \cdot \vartheta}{\vartheta + 243{,}12}} \quad [\text{Pa}] \\[4ex]
e(\vartheta, rF) &= \frac{rF}{100} \cdot e_{\text{sat}}(\vartheta) \quad [\text{Pa}]
\end{align}
folgt daraus die kompakte Form der absoluten Feuchte unter Berücksichtigung des Drucks:
\begin{align}
\text{aF}(\vartheta, rF, p) &= 622 \cdot \frac{e(\vartheta, rF)}{p - e(\vartheta, rF)} \quad \left[\frac{\text{g}}{\text{kg}}\right]
\end{align}
wobei
\begin{align*}
\vartheta &= \text{Temperatur in } ^\circ\text{C} \\[3ex]
rF &= \text{Relative Luftfeuchtigkeit in } \% \, (0 \dots 100) \\[3ex]
p &= \text{Tatsächlicher Gesamtdruck (Betriebsdruck) in Pa} \\[3ex]
e_{\text{sat}} &= \text{Sättigungsdampfdruck in Pa} \\[3ex]
e &= \text{Tatsächlicher Wasserdampfpartialdruck in Pa}
\end{align*}


\section{Dämmschichttrockung}
Der Druck in der Prozessluft ist beim Saugen geringer als der Durck im Raum. Der Sättigungsdampfdruck von Wasser beträgt bei einer normalen Raumtemperatur von 20 °C ca. 2.340 Pascal. Damit kann das Wasser verdampfen. 

Der Gesamtdruck in der Dämmschicht sinkt durch den Sauger leicht ab (z. B. von 1013 mbar auf ca. 950 mbar absolut).

\textbf{Frage} Haben wir eine Chance, den Druck am Sensor in der Prozessluft zu messen?

\begin{itemize}
    \item \textbf{Der Effekt auf die Feuchte:} Weil der Druck leicht sinkt, dehnt sich die Luft aus. Ihre Dichte sinkt ein wenig (sie wird geringer als die 1,204 kg/m³).
    \item \textbf{Das Ergebnis in der Praxis:} Die Luft kann durch den veränderten Druck pro Kubikmeter sogar geringfügig mehr Wasserdampf aufnehmen, und die Verdunstung des Wassers in der Dämmschicht wird beschleunigt.
    \item \textbf{Für die Messung:} Die unkorrigierte Magnus-Formel liefert  immer noch einen praxistauglichen Näherungswert. 
\end{itemize}

\textbf{ Diese Aussage sollte hinterfragt werden.}

\section{Ampelstatus}

\subsection{Ampel Report}
Ampelreport (prüfen!!):

\begin{table}[htbp]
    \centering
    \caption{Übersicht der der Ampelstatus} % Optional: Tabellenüberschrift
    \label{tab:ampelstatus}                  % Optional: Label für Verweise
    \begin{tabular}{lll}
        \toprule
        Status & relative Feuchte & Differenz absolute Feuchte \\
        \midrule
        Rot  & $\text{rF}_i \ge 65\,\% $       & $\Delta\text{aF} > 3 \frac{\text{g}}{\text{m}^3}$ \\ \addlinespace 
        Gelb & $65 > \text{rF}_i > 50$         & $\Delta\text{aF} \le 3\,\frac{\text{g}}{\text{m}^3}$ \\ \addlinespace
        Grün & $50 \ge \text{rF}_i$            & $\Delta\text{aF} \le 2\,\frac{\text{g}}{\text{m}^3}$ \\ \addlinespace
        nicht definiert & sonst                           & sonst \\  \bottomrule
    \end{tabular}
\end{table}


\subsection{Berechnung des Ampelstatus}
Bei der Berechnung des Ampelstatus wird der Absolutwert der Differenz der absoluten Feuchten in der Prozessluft $\text{aF}_p$ (process) und in der Raumluft $\text{aF}_a$ (ambient) gebildet.

\begin{align}
    \Delta\text{aF} 
    &=  \left| \text{aF}_p - \text{aF}_a \right| 
    \\[4ex]
    &=
    \frac{\text{rF}_p \cdot 611{,}2\, \frac{\text{J}}{\text{m}^3} \cdot e^{\dfrac{17{,}62 \cdot \vartheta_p}{\vartheta_p + 243{,}12}}}
    {461{,}52 \frac{\text{J}}{\text{kg} \cdot \text{K}} \cdot (\vartheta_p + 273{,}15)\, \text{K}}
    \cdot \frac{10}{1{,}204\, \frac{\text{kg}}{\text{m}^3}}
    - 
    \frac{\text{rF}_a \cdot 611{,}2\, \frac{\text{J}}{\text{m}^3} \cdot e^{\dfrac{17{,}62 \cdot \vartheta_a}{\vartheta_a + 243{,}12}}}
    {461{,}52 \frac{\text{J}}{\text{kg} \cdot \text{K}} \cdot (\vartheta_a + 273{,}15)\, \text{K}}
    \cdot \frac{10}{1{,}204\, \frac{\text{kg}}{\text{m}^3}}
    \\[4ex]
    &=
    \left[
    \text{rF}_p  \cdot
    \frac{ e^{\dfrac{17{,}62 \cdot \vartheta_p}{\vartheta_p + 243{,}12}}}
    { (\vartheta_p + 273{,}15)\, \text{K}}
    - 
    \text{rF}_a \cdot
    \frac{ e^{\dfrac{17{,}62 \cdot \vartheta_a}{\vartheta_a + 243{,}12}}}
    {(\vartheta_a + 273{,}15)\, \text{K}}
    \right]
    \cdot \frac{611{,}2\, \frac{\text{J}}{\text{m}^3}}{461{,}52 \frac{\text{J}}{\text{kg} \cdot \text{K}}}
    \cdot \frac{10}{1{,}204\, \frac{\text{kg}}{\text{m}^3}}
    \\[4ex]
     &=
    \left[
    \frac{\text{rF}_p}{\vartheta_a + 273{,}15} \cdot
    e^{f(\vartheta_p)}
    -
    \frac{\text{rF}_a}{\vartheta_p + 273{,}15} \cdot
    e^{f(\vartheta_a)}
    \right]
    \cdot \text{Konstante}   
\end{align}

\subsection{Graphische Darstellung}
\subsubsection{Ansatz 1}
2D Plot der Funktion im Intervall $[0 \dots 15]$:
\begin{equation}
    \Delta\text{aF} 
    =  \left| \text{aF}_p - \text{aF}_a \right|     
\end{equation}



\newpage
% Hier wird das Literaturverzeichnis gedruckt:
\printbibliography[heading=bibintoc, title={Literaturverzeichnis}]

\end{document}